\subsection{High-object reductions at sixteen arrows}

\begin{lemma}[A singleton tree with one nontrivial corner]
\label{C:tree-singleton-core-cups}
Suppose a connected finite category has at least two objects,
its undirected cross-arrow support is a tree, every nonempty
cross-hom set is a singleton, and at most one corner monoid
is nontrivial.  Every positive-degree rational product on its
classifying space is zero.
\end{lemma}
\begin{proof}
If an edge has arrows in both directions, any edge connecting
one of its vertices to a third vertex would force the other
side of a triangle by composition.  Connectedness of the tree
therefore leaves only those two objects.  At least one corner
is trivial, so that object is both initial and terminal and
the category is contractible.  We may exclude this case.

A directed path through three distinct objects would force
its composite cross-hom set and hence a triangle in the
undirected support.  No such path occurs.  Thus the nerve
is the graph of correspondence spaces indexed by that tree.
Remove a leaf whose corner and neighboring corner are both
trivial; its correspondence is a point and its mapping
cylinder is an interval, so this is a homotopy-equivalence
collapse.  Repeating leaves a star around the possible
nontrivial corner $M$, with $d\geq1$ trivial leaves.
Each edge correspondence is $BM$ and its map to the center
is the identity, because its unique cross arrow is fixed
by all actors.  The resulting space is the join of $BM$
with a discrete $d$-point set.  It is contractible for $d=1$
and a wedge of $d-1$ suspensions of $BM$ for $d>1$.
Its positive-degree products vanish.  If every corner is
trivial, the original tree itself is contractible.
\end{proof}

\begin{lemma}[Semisimple diagonal choices with many objects]
\label{C:sixteen-high-semisimple}
A finite category with sixteen arrows, at least five objects,
acyclic support, and semisimple rational corner algebras
cannot have classifying space $\mathbb CP^2$.
\end{lemma}
\begin{proof}
The relative normalized bar over the product of the corner
algebras is an exact rational cohomology model.  Its length
is at most the height of the support order.  A four-step
chain requires ten distinct cross-hom sets.  With $k$ objects,
connecting the other $k-5$ objects requires at least $k-5$
more cross arrows; together with identities this costs
$2k+5$.  For $k\geq6$ it exceeds sixteen, so the relative
bar has length at most three and fourth homology vanishes.

For $k=5$, a four-step chain forces all ten forward hom-sets.
If all corners are trivial, Theorem~\ref{C:all-acyclic21}
applies.  Otherwise exactly one corner can have order two,
all others are trivial, and every cross-hom set is a singleton.
At least one extreme object of the chain has trivial corner;
it is respectively initial or terminal.  The category is
contractible.  These cases exhaust the sixteen-arrow budget.
\end{proof}

\begin{theorem}[Seven or more objects at sixteen arrows]
\label{C:sixteen-seven-objects}
Every finite category with at least seven objects and at most
sixteen arrows fails to have classifying space $\mathbb CP^2$.
\end{theorem}
\begin{proof}
Smaller arrow totals are already excluded.  Assume connectedness.
The tree budget leaves only seven or eight objects.  For $k$
objects, write $e=\sum_c(|\operatorname{End}(c)|-1)$ for the
number of nonidentity corner arrows.  Relative to the minimum
$2k-1$ arrows, there are $17-2k$ extra arrows.

We first dispose of nontrivial strongly connected components.
A retract object can be deleted, reducing to the already
proved lower bound.  Without a retract, every corner in a
nontrivial strong component has order at least two.  A
component of $r$ objects also has all $r(r-1)$ directed cross
hom-sets nonempty.  Compared with a connecting tree on those
objects, these requirements cost at least
\[
 r+\bigl(r(r-1)-(r-1)\bigr)=r+(r-1)^2
\]
extra arrows.  This is at least seven for $r\geq3$ and three
for $r=2$.  At eight objects there is only one extra arrow,
so no such component remains.  At seven objects there are
three extras: a possible strong pair has both corner orders
two, one arrow each way, and every other corner and cross
hom-set is trivial or a singleton.  If an opposite composite
is a unit, it gives a retract.  Otherwise each corner is the
identity-adjoined zero monoid.  Its zero idempotent fixes every
incoming cross arrow, since each such hom-set is a singleton.
Lemma~\ref{C:one-sided-corner-deletion} removes a vertex of
the pair, again reducing the arrow count.  We may therefore
assume acyclic support.

For eight objects, $e\leq1$, so every corner has order at
most two and every corner algebra is semisimple over
$\mathbb Q$.  Lemma~\ref{C:sixteen-high-semisimple} applies.
For seven objects, $e\leq3$.  If all corner orders are at
most two, the same lemma applies.  Thus only a corner of
order three or four needs further attention.

If $e=3$, there are exactly six cross arrows.  Connectedness
makes their undirected support a tree and each cross-hom a
singleton.  Every corner has order at most four and hence
rationally acyclic classifying space, by the complete small
corner classification.  Every edge correspondence in the
graph-of-spaces description is the product of the two adjacent
corner classifying spaces, and is likewise rationally acyclic.
The Mayer--Vietoris sequence along the tree gives rational
acyclicity of the whole category.

The remaining case has $e=2$, one order-three corner, and
all other corners trivial.  There are seven cross arrows.
If they are in distinct hom-sets, projection to the support
poset is a rational homology equivalence.  Indeed, filter the
nerve by its strict object chain: over each such chain the
corner bars are products of rationally acyclic classifying
spaces, and the singleton cross arrows give their augmentation
maps.  The associated homology spectral sequence is precisely
the ordinary support-poset complex.  A four-object chain would
already require six cross arrows and three further connecting
arrows, so this support has height at most two.  Its fourth
homology is zero.

Otherwise there is exactly one two-element cross-hom set,
all others are singletons, and the undirected support is a
tree.  It has no directed path through three distinct objects.
All edge correspondences except possibly the doubled one are
rationally acyclic.  The doubled correspondence is either a
discrete two-point space or the action category of a two-state
act over the order-three corner.  Its rational homology is
zero in degrees at least two.  To check this last assertion,
the complete order-three corner list consists of group-type
monoids, whose trivial rational modules are projective, and
the identity-adjoined left- and right-zero bands of order two.
The latter rational algebras are triangular two-by-two matrix
algebras and have global dimension one.  The two-sided action
bar therefore has no homology in degrees at least two in
every case.  Mayer--Vietoris for the graph of correspondences
now gives zero fourth homology of the whole category.  This
exhausts $e=2,3$, and proves the theorem.
\end{proof}
